% shabbat-siddur/sections/front-matter.tex
% Title page (recto) and copyright page with genizah notice (verso).

% ---- Title page ------------------------------------------------------------
\thispagestyle{empty}
\vspace*{\fill}
\begin{center}
  {\Large\bfseries תְּפִלּוֹת מִנְחָה וְעַרְבִית\\ לְמוֹצָאֵי שַׁבָּת\par}
  \medskip
  \begin{otherlanguage}{english}
    {\large\itshape Shabbat Mincha, Arvit\\ and Havdalah\par}
    \bigskip
    {\small According to the custom of the\\
     Turkish Sephardic community of Seattle\par}
  \end{otherlanguage}
\end{center}
\vspace*{\fill}
\begin{center}
  \begin{otherlanguage}{english}
    {\small\scshape Sephardic Liturgical Press\par}
  \end{otherlanguage}
\end{center}
\clearpage

% ---- Copyright page --------------------------------------------------------
\thispagestyle{empty}
\begin{otherlanguage}{english}
\footnotesize\raggedright
\vspace*{\fill}

\textit{Tefillot Mincha VeArvit LeMotzei Shabbat}\\
According to the custom of the Turkish Sephardic community of Seattle

\medskip
© 2026 Sephardic Liturgical Press. All rights reserved in the typesetting,
layout, and English rubrics of this edition.

\medskip
The Hebrew text is based on \textit{Siddur Edot HaMizrach} (Torat Emet), a
public-domain text made available by Sefaria, and has been edited to follow
the Turkish Sephardic custom of Seattle. Psalm 144 is from the public-domain
Westminster Leningrad Codex text (tanach.us).

\medskip
Set in \pressHebrewFontName{} (Culmus Project) and Gentium Plus
(SIL International).

\medskip
First edition, 5787 / 2026.

\vspace*{\fill}
\end{otherlanguage}

\begin{center}
  \fbox{\parbox{0.85\linewidth}{\centering
    סֵפֶר זֶה מֵכִיל שֵׁמוֹת קְדוֹשִׁים.\\ נָא לִשְׁמֹר עַל קְדֻשָּׁתוֹ.\par
    \smallskip
    \begin{otherlanguage}{english}\footnotesize\centering
      This siddur contains the Holy Name. Please treat it with respect,
      and when it is worn, place it in a genizah rather than discarding it.\par
    \end{otherlanguage}}}
\end{center}
\clearpage
